% =====================================================================================
% Chapter 5 - concise version
% Required packages in the preamble: graphicx, booktabs, amsmath (already standard).
% Figure placeholders: replace each \fbox{...} with \includegraphics[...]{<path>} later.
% \tbd marks values to be filled in after the benchmark is run with all seven checkpoints.
% =====================================================================================
\providecommand{\tbd}{\underline{\hspace{1.2cm}}}

\chapter{System Design and Implementation}
\label{chap:system_design}

This chapter describes how the RB-Net framework presented in Chapter~4 is realized as an interactive web-based system for missense variant pathogenicity classification. It covers the overall system architecture, the integration of foundation models and bioinformatics resources into a unified inference engine, the data flow and communication protocols that support single-variant and batch analysis, and the data management and deployment strategy. The chapter concludes with an overview of the web interface and an empirical evaluation of the system's runtime latency.

% -------------------------------------------------------------------------------------
\section{System Architecture}
\label{sec:system_architecture}

The system follows a three-tier client--server architecture, as illustrated in Fig.~\ref{fig:system_architecture}, which separates user interaction, AI orchestration, and data management so that computationally intensive annotation and inference do not block the user interface. The presentation tier is a React single-page application organized into four views: a system overview, single-variant inspection, batch screening, and explainability analysis. It communicates with the backend through REST requests for single-variant queries and Server-Sent Events (SSE) for streaming batch results, while variant and feature configurations are kept in the URL so that each analysis can be shared and revisited.

The application tier is a FastAPI service that orchestrates the complete inference workflow, from variant annotation and foundation model feature extraction to checkpoint routing and Shapley-based explanation. Blocking operations, such as external annotation processes and GPU computation, are executed in background worker threads to keep the service responsive under concurrent requests. The data tier combines a PostgreSQL database extended with \texttt{pgvector}, which persists variant annotations, embeddings, and cached predictions, with locally deployed bioinformatics resources: an offline Ensembl Variant Effect Predictor (VEP) service, the GRCh38 reference genome and proteome, precomputed SpliceAI scores, and FAISS reference indexes. Keeping all resources local ensures reproducible annotation and removes any dependence on external web services.

\begin{figure}[htbp]
    \centering
    % \includegraphics[width=0.95\textwidth]{figures/ch5/fig5_1_system_architecture.pdf}
    \fbox{\parbox{0.92\textwidth}{\centering\vspace{2.2cm}\textit{[Figure 5.1: Three-tier system architecture]}\vspace{2.2cm}}}
    \caption{High-level three-tier architecture of the RB-Net system, comprising the React presentation tier, the FastAPI inference tier, and the PostgreSQL and bioinformatics data tier.}
    \label{fig:system_architecture}
\end{figure}

% -------------------------------------------------------------------------------------
\section{Artificial Intelligence Model and Bioinformatics Integration}
\label{sec:ai_integration}

\subsection{Feature Extraction Pipeline}

To avoid cold-start latency, the backend loads the foundation models, sequence indexes, FAISS indexes, and normalization artifacts once at startup and performs a warm-up forward pass before accepting requests. Each incoming variant, specified as \texttt{chr:pos:ref>alt} on GRCh38, is then annotated by the offline VEP service, which selects the MANE Select or canonical transcript and provides the consequence, the affected protein position, and the eleven core biological features: two population allele frequencies from the VEP cache and nine evolutionary conservation scores from dbNSFP. AlphaMissense, REVEL, and ClinVar annotations are retrieved at the same time for benchmarking.

Since RB-Net is trained exclusively on missense single-nucleotide variants, a missense gatekeeper rejects every other consequence type before inference, and the reference allele and amino acid are verified against the local GRCh38 genome and proteome to detect inputs from an incorrect genome build. For accepted variants, the nine SpliceAI features are read directly from the compressed SpliceAI archive through a custom tabix reader that decompresses only the blocks surrounding the variant. The DNA and protein windows ($1{,}001$~bp and $501$~aa) centered on the mutation are then encoded by NT-v1 500M\footnote{\url{https://huggingface.co/InstaDeepAI/nucleotide-transformer-500m-human-ref}} and ESM-2 650M\footnote{\url{https://huggingface.co/facebook/esm2_t33_650M_UR50D}} in FP32 precision, and the resulting reference and alternative embeddings are used to compute the geometric descriptors, including the local intrinsic dimensionality estimated from $k = 64$ nearest neighbors in the FAISS indexes.

\subsection{Ablation-Aware Model Routing}

The web interface exposes three feature groups, each controlled by a single switch: \textit{DNA}, \textit{Protein}, and \textit{Bio + Geom + SpliceAI}. Each combination is identified by a three-bit key, and since disabling all groups is not allowed, the system supports seven configurations, summarized in Table~\ref{tab:ablation_configs}. A checkpoint registry maps each key to its trained checkpoint and infers the model architecture directly from the stored weights, so that new checkpoints can be activated without code changes. Internally, the key is stored as a four-bit mask in which the Bio and Geometry bits are always set together.

The full configuration uses the \texttt{RouterExpertCorrection} model ($6{,}462{,}728$ parameters, router temperature $T = 0.3074$, learnable correction weight $\beta$), while the six reduced configurations use \texttt{MultiStrategyFusionModel} checkpoints. Each checkpoint applies its own MCC-optimal decision threshold $\tau$ (e.g., $\tau = 0.6376$ for the full configuration), and a variant is classified as \textit{Pathogenic} if its predicted probability satisfies $p \geq \tau$.

\begin{table}[htbp]
    \centering
    \caption{Feature-ablation configurations supported by the checkpoint registry.}
    \label{tab:ablation_configs}
    \begin{tabular}{clcc}
        \toprule
        \textbf{Key} & \textbf{Active feature groups} & \textbf{Architecture} & \textbf{SHAP players} \\
        \midrule
        \texttt{111} & DNA, Protein, Bio + Geom + SpliceAI & RouterExpertCorrection & 30 \\
        \texttt{110} & DNA, Protein & MultiStrategyFusion & 2 \\
        \texttt{101} & DNA, Bio + Geom + SpliceAI & MultiStrategyFusion & 29 \\
        \texttt{011} & Protein, Bio + Geom + SpliceAI & MultiStrategyFusion & 29 \\
        \texttt{100} & DNA & MultiStrategyFusion & 1 \\
        \texttt{010} & Protein & MultiStrategyFusion & 1 \\
        \texttt{001} & Bio + Geom + SpliceAI & MultiStrategyFusion & 28 \\
        \bottomrule
    \end{tabular}
\end{table}

\subsection{Integrated Explainability}

Instead of relying on an external explainability library, the inference engine embeds a permutation-based Shapley value estimator that runs directly on the GPU. Each active embedding displacement ($\boldsymbol{\delta}_{\text{DNA}}$ or $\boldsymbol{\delta}_{\text{Prot}}$) is treated as a single player, and each of the $28$ standardized tabular features is treated as an individual player, giving $n = 30$ players for the full configuration. Using a zero baseline, the engine samples $128$ antithetic permutations, evaluates all $128 \times (n + 1)$ nested coalitions as one batched tensor ($3{,}968$ forward passes for the full configuration), and returns logit-scale attributions that satisfy the efficiency property
\[
    \sum_{i=1}^{n} \phi_i = f(\mathbf{x}) - f(\mathbf{0}).
\]
Because the attributions are computed in the same request as the prediction, every result shown to the user is accompanied by its feature-level explanation.

% -------------------------------------------------------------------------------------
\section{Data Flow and Communication Protocols}
\label{sec:data_flow}

\subsection{Hybrid Lazy-Caching}

Interactive exploration requires fast responses, whereas analyzing a new variant involves costly annotation and foundation model inference. The system therefore caches upstream features and downstream predictions separately, and each request follows one of three paths. In a \textit{full cache hit}, a prediction for the same variant, configuration, and model version already exists and is returned directly from the database. In a \textit{feature reuse} path, triggered when the user changes the feature configuration or enables \textit{Live Bypass} mode, the stored embeddings and raw features are reloaded and only normalization, classification, and explanation are re-executed, without repeating VEP annotation or foundation model inference. In a \textit{full cache miss}, the complete pipeline is executed and its outputs are persisted for later requests.

Cached entries are versioned to prevent stale results. Feature records carry a fingerprint of the extraction settings (model identifiers, window sizes, device, and precision), and prediction records additionally carry the checkpoint identity, so that any change to the extraction pipeline or to a model checkpoint automatically invalidates the affected entries.

\subsection{Single-Variant Inference Workflow}

The single-variant workflow, shown in Fig.~\ref{fig:single_sequence}, is designed to overlap slow external annotation with independent GPU work. The user submits a variant and a feature configuration at stage~1, and the backend checks the cache at stage~2, returning immediately on a hit. On a miss, the VEP annotation is dispatched to a background thread at stage~3, while the DNA window extraction, SpliceAI lookup, and NT-v1 inference run concurrently at stage~4. Once annotation completes, the missense gatekeeper validates the variant at stage~5, after which the protein window is extracted and encoded by ESM-2 at stage~6. The geometric descriptors are computed and all upstream features are persisted at stage~7, the routed checkpoint produces the prediction and its Shapley attributions at stage~8, and the result is stored and returned to the browser as a JSON response at stage~9.

\begin{figure}[htbp]
    \centering
    % \includegraphics[width=0.95\textwidth]{figures/ch5/fig5_2_single_variant_workflow.pdf}
    % Stages to draw: (1) submit variant + configuration, (2) cache lookup (hit -> return),
    % (3) VEP in background thread, (4) DNA window + SpliceAI + NT-v1 in parallel,
    % (5) missense gatekeeper, (6) protein window + ESM-2, (7) geometry + persist features,
    % (8) classifier + SHAP, (9) persist prediction + JSON response.
    \fbox{\parbox{0.92\textwidth}{\centering\vspace{2.2cm}\textit{[Figure 5.2: Single-variant inference workflow (stages 1--9)]}\vspace{2.2cm}}}
    \caption{Single-variant inference workflow with hybrid lazy-caching, in which VEP annotation overlaps with DNA feature extraction.}
    \label{fig:single_sequence}
\end{figure}

\subsection{Batch Screening via Server-Sent Events}

Batch screening streams results incrementally through SSE, which avoids HTTP timeouts and lets the user monitor progress in real time; the workflow is shown in Fig.~\ref{fig:batch_sse}. At stage~1, the user uploads a variant file (VCF, CSV, TXT, Parquet, or Excel) or selects the built-in $10$-variant benchmark cohort, and at stage~2 the backend extracts and deduplicates the variant coordinates. At stage~3, all uncached variants are annotated in a single VEP run, followed by the missense gatekeeper and a batched SpliceAI lookup. The accepted variants are then processed in micro-batches of up to eight variants, which are encoded jointly by both foundation models at stage~4, after which their geometric descriptors are computed and persisted in a single database transaction at stage~5. At stage~6, each variant is classified and explained, and at stage~7 its result is emitted immediately as an SSE event, while rejected variants are reported as skipped events with the reason for rejection. The browser updates the live triage table and saves the progress in the session at stage~8, and a final summary event closes the stream at stage~9.

\begin{figure}[htbp]
    \centering
    % \includegraphics[width=0.95\textwidth]{figures/ch5/fig5_3_batch_sse_workflow.pdf}
    % Stages to draw: (1) upload file / demo cohort, (2) parse + deduplicate,
    % (3) cohort-wide VEP + gatekeeper + SpliceAI, (4) micro-batched NT-v1 / ESM-2 (N <= 8),
    % (5) geometry + persist micro-batch, (6) per-variant classifier + SHAP,
    % (7) SSE event (variant_stream / variant_skipped), (8) live triage table + session storage,
    % (9) stream_completed summary.
    \fbox{\parbox{0.92\textwidth}{\centering\vspace{2.2cm}\textit{[Figure 5.3: Batch screening and SSE streaming workflow (stages 1--9)]}\vspace{2.2cm}}}
    \caption{Batch screening workflow combining cohort-wide annotation, micro-batched foundation model inference, and incremental result streaming through SSE.}
    \label{fig:batch_sse}
\end{figure}

% -------------------------------------------------------------------------------------
\section{Data Management and Deployment}
\label{sec:db_deployment}

\subsection{Database Design}

The database, shown in Fig.~\ref{fig:db_erd}, consists of three related tables. The \texttt{variants} table stores each GRCh38 variant together with its gene, consequence, protein change, and normalized ClinVar label. The \texttt{inference\_features} table keeps a one-to-one record per variant containing the four $1{,}280$-dimensional reference and alternative embeddings as \texttt{pgvector} columns, the unscaled biological and geometric features, and the extraction fingerprint. Because the classifiers were trained on embeddings stored in FP16, the extracted embeddings are rounded to FP16 precision before being stored as FP32 vectors, keeping online features consistent with the training data. The \texttt{predictions} table holds at most one record per variant and feature configuration, containing the predicted probability, label, benchmark scores, Shapley attributions, and model version. All writes are performed as atomic upserts, and upstream features are committed before classification, so that a failure in a later stage never forces the costly annotation and embedding steps to be repeated.

\begin{figure}[htbp]
    \centering
    % \includegraphics[width=0.85\textwidth]{figures/ch5/fig5_4_database_erd.pdf}
    \fbox{\parbox{0.92\textwidth}{\centering\vspace{2.2cm}\textit{[Figure 5.4: Entity-relationship diagram of the RB-Net database]}\vspace{2.2cm}}}
    \caption{Entity-relationship diagram of the RB-Net database, showing the one-to-one relationship between \texttt{variants} and \texttt{inference\_features} and the one-to-many relationship with \texttt{predictions}.}
    \label{fig:db_erd}
\end{figure}

\subsection{Deployment Strategy}

The system adopts a hybrid deployment that containerizes infrastructure services while running GPU inference natively. PostgreSQL and the offline VEP service run in Docker containers, which isolates their dependencies and allows the reference data to be mounted from the host, whereas the FastAPI backend runs in a local Python environment to obtain direct, low-overhead access to the CUDA runtime. The backend is started without hot reloading so that the multi-gigabyte foundation models are loaded on the GPU only once, and a single startup script launches the database, backend, and frontend in sequence. The main technologies are summarized in Table~\ref{tab:tech_stack}.

\begin{table}[htbp]
    \centering
    \caption{Technology stack of the RB-Net system.}
    \label{tab:tech_stack}
    \begin{tabular}{ll}
        \toprule
        \textbf{Layer} & \textbf{Technologies} \\
        \midrule
        Frontend & React 19, TypeScript, Vite, Tailwind CSS, Recharts \\
        Backend & Python 3.10, FastAPI, Uvicorn, Pydantic, SQLAlchemy \\
        Deep learning & PyTorch 2.4 (CUDA 12.1), Hugging Face Transformers, FAISS \\
        Database & PostgreSQL 16 with \texttt{pgvector} (Docker) \\
        Bioinformatics & Ensembl VEP 116 with dbNSFP 5.3.1a (Docker), SpliceAI MANE v1.4, pyfaidx \\
        \bottomrule
    \end{tabular}
\end{table}

% -------------------------------------------------------------------------------------
\section{Web Interface}
\label{sec:web_interface}

The web interface, shown in Fig.~\ref{fig:web_interface}, turns the inference engine into four task-oriented views. The overview page in Fig.~\ref{fig:web_interface}a introduces the system and links to each module. The single-variant view in Fig.~\ref{fig:web_interface}b accepts a GRCh38 coordinate together with the three feature switches and displays the predicted probability and label, the ClinVar label for reference, and a comparison with AlphaMissense and REVEL; a \textit{Live Bypass} option forces recomputation instead of reading the cache. The batch view in Fig.~\ref{fig:web_interface}c accepts a variant file, shows the streaming progress, the number of accepted and rejected variants, and a live triage table ranking the most likely pathogenic variants. Finally, the explainability view in Fig.~\ref{fig:web_interface}d presents the Shapley attributions of a selected variant as a diverging bar chart, separating features that increase or decrease the predicted pathogenicity. Each result can be opened from one view to another with its variant and feature configuration preserved.

\begin{figure}[htbp]
    \centering
    % Replace each box with \includegraphics[width=\linewidth]{figures/ch5/fig5_5x_....png}
    \begin{minipage}{0.48\textwidth}\centering
        \fbox{\parbox{0.95\linewidth}{\centering\vspace{1.6cm}\textit{[Overview page]}\vspace{1.6cm}}}\\[2pt]
        \small (a) System overview
    \end{minipage}\hfill
    \begin{minipage}{0.48\textwidth}\centering
        \fbox{\parbox{0.95\linewidth}{\centering\vspace{1.6cm}\textit{[Single-variant view]}\vspace{1.6cm}}}\\[2pt]
        \small (b) Single-variant inspection
    \end{minipage}\\[8pt]
    \begin{minipage}{0.48\textwidth}\centering
        \fbox{\parbox{0.95\linewidth}{\centering\vspace{1.6cm}\textit{[Batch view]}\vspace{1.6cm}}}\\[2pt]
        \small (c) Batch screening
    \end{minipage}\hfill
    \begin{minipage}{0.48\textwidth}\centering
        \fbox{\parbox{0.95\linewidth}{\centering\vspace{1.6cm}\textit{[Explainability view]}\vspace{1.6cm}}}\\[2pt]
        \small (d) Explainability analysis
    \end{minipage}
    \caption{Main views of the RB-Net web interface.}
    \label{fig:web_interface}
\end{figure}

% -------------------------------------------------------------------------------------
\section{Latency Evaluation}
\label{sec:latency}

Runtime latency was measured on a workstation with an AMD Ryzen~7 8745H CPU, 32~GB of RAM, and an NVIDIA GeForce RTX 4060 Laptop GPU (8~GB VRAM) running Windows~11, with all inference performed in FP32. Startup preloading takes approximately \tbd{}~s and occupies approximately \tbd{}~GB of GPU memory, and is excluded from the measurements. % TODO: fill in after the FP32 benchmark
Latency was recorded by the backend's built-in profiler, from the moment a validated request enters the pipeline until its response is constructed, under a sequential single-user workload on the $10$-variant benchmark cohort. Network transfer and browser rendering are therefore not included. Table~\ref{tab:latency_benchmarks} reports the results for each execution path.

% TODO: fill in the latency column once all seven checkpoints are available.
\begin{table}[htbp]
    \centering
    \caption{Server-side processing latency across execution paths on the reference workstation.}
    \label{tab:latency_benchmarks}
    \begin{tabular}{llr}
        \toprule
        \textbf{Scenario} & \textbf{Execution path} & \textbf{Latency} \\
        \midrule
        Single new variant & Full cache miss & \tbd{} s / variant \\
        Batch cohort ($10$ variants) & Full cache miss & \tbd{} s total \\
        Configuration change / Live Bypass & Feature reuse & \tbd{} s / query \\
        Repeated query & Full cache hit & \tbd{} ms / query \\
        \bottomrule
    \end{tabular}
\end{table}

The results show the benefit of separating feature extraction from classification. Batch processing reduces the average cost to \tbd{}~s per variant, a \tbd{}$\times$ speed-up over sequential single-variant requests, mainly because VEP is invoked once for the whole cohort and the foundation models process several variants per forward pass. Once a variant's features are stored, changing the feature configuration takes only \tbd{}~s, and repeated queries are answered in \tbd{}~ms, enabling responsive interactive exploration and model interpretation.
